% GENERATED FILE. Edit canonical lessons, then run npm run generate:books.
% canonical-source-hash: fnv1a64:c46235160b7abe76
% canonical-lessons: ZH-C89-gei, ZH-W89-gei, ZH-C89-zuo, ZH-C89-fang, ZH-C89-dai, ZH-C89-zhan, ZH-C89-practice

\chapter{Give Me That}
\label{ch:zh-give-me-that}
\hlchaptermodality{89}

\begin{chapteropening}
\textbf{By the end of this chapter:} \emph{I can name everyday actions: give, do, put, bring, stand.}

Five everyday actions, give, do, put, bring and stand, and \zh{给我} as a request, built on one new written character: \zh{给}.
\end{chapteropening}

\section[\texorpdfstring{\zh{给} (gěi)}{给 (gěi)}]{\zh{给} (gěi) — to give; to, for}

\label{lesson:ZH-C89-gei}



\begin{quote}
A new chapter, and it opens with the verb that hands things over.

\begin{itemize}
\raggedright
  \item \emph{Recall:} say \zh{两个人} — \emph{liǎng ge rén}, two people
  \item \emph{Recall:} say \emph{liǎng zhī māo}, two cats — \textbf{R1}, the chapter you have finished
  \item \emph{Recall:} say \emph{sòng}, to give as a present — \textbf{R3}, durable
\end{itemize}
\end{quote}

\subsection*{You'll want to know: \zh{给}}
\textbf{\zh{给}} — \emph{gěi} — \textquotedblleft{}to give; to, for\textquotedblright{}.

\begin{quote}
\textbf{\zh{妈妈给我一个手机}} \emph{māma gěi wǒ yí ge shǒujī} — Mum gives me a phone.
\end{quote}

\begin{quote}
\textbf{\zh{给我}} \emph{gěi wǒ} — Give it to me.
\end{quote}

The next lesson puts the new character in your hand.

\begin{cousinweb}
\textbf{\zh{给}} (\emph{gěi}) is in the Guangyun as \emph{kip} (Guangyun, in Baxter's transcription), glossed 'to supply'. That reading regularly gives jǐ, still heard in a few written words; how the everyday gěi came about is not settled.
\end{cousinweb}

\begin{grammarlens}[title={What you've built}]
The person comes first, then the thing: \zh{给}, \zh{我}, \zh{一个手机}. On its own, \zh{给我} is a request: give me. \emph{sòng} gives a present; \zh{给} gives anything at all.
\end{grammarlens}

\subsection*{Your turn}
\begin{itemize}
\raggedright
  \item \emph{Say it:} \emph{gěi}
  \item \emph{Say it:} \emph{gěi}, once more
  \item \emph{Say it:} \zh{妈妈给我一个手机} — \emph{māma gěi wǒ yí ge shǒujī} — Mum gives me a phone.
  \item \emph{Say it:} \zh{给我} — \emph{gěi wǒ} — Give it to me.
  \item \emph{Recall:} say \emph{sòng}, then \emph{gěi} again
\end{itemize}

\begin{tcolorbox}[breakable,colback=teal!4,colframe=teal!35!black,title={Before you move on}]
What does \zh{给} mean? (\textquotedblleft{}To give; to, for\textquotedblright{}.) Which comes first after \zh{给}, the person or the thing? (\textbf{The person}.)
\end{tcolorbox}
\section[\texorpdfstring{\zh{给} (gěi)}{给 (gěi)}]{\zh{给} — the silk piece, then \zh{合}}

\label{lesson:ZH-W89-gei}



\begin{quote}
The word first, then the character inside it.

\begin{itemize}
\raggedright
  \item \emph{Recall:} say \emph{gěi} — \textbf{R1}, one lesson back, before it decays
\end{itemize}

\begin{quote}
\zh{人} \zh{一} \zh{口}
\end{quote}

\emph{Read it:} each one aloud, with nothing said first — \emph{rén}, \emph{yī}, \emph{kǒu}
\end{quote}

\subsection*{Script — the silk piece, then \zh{合}}
\hlblockfigure{\detokenize{figures/ZH-W89-gei-filmstrip.pdf}}{How \zh{给} is written, stroke by stroke}

\begin{quote}
\zh{给}
\end{quote}

\textbf{Two pieces, left before right. The left piece is the one new shape.}

\textbf{The silk piece \zh{纟}, three strokes}: a stroke that falls to the left and turns right, a second, longer one the same way, then a stroke rising to the right.

\textbf{\zh{合}, six strokes}: a roof like \zh{人}, a left-falling stroke and a right-falling one; a short horizontal under it; then \zh{口}.

3 + 6 = nine.

\subsection*{Your turn}
\begin{itemize}
\raggedright
  \item \emph{Write it:} \zh{给}, all nine, with the model in front of you
  \item \emph{Cover:} write it again from memory, and count the pieces, not the strokes
\end{itemize}

\begin{tcolorbox}[breakable,colback=teal!4,colframe=teal!35!black,title={Before you move on}]
How many strokes in \zh{给}? (\textbf{Nine}.) Which piece is new? (\textbf{The silk piece, on the left}.) What sits at the foot of the right piece? (\textbf{\zh{口}}.)

Source: \href{https://raw.githubusercontent.com/chanind/hanzi-writer-data/68d10a4b21150cae5e1ebbd223eed289cf32d90c/data/\%E7\%BB\%99.json}{Hanzi Writer Data \zh{给}, pinned stroke paths and medians}.
\end{tcolorbox}
\section[\texorpdfstring{\zh{做} (zuò)}{做 (zuò)}]{\zh{做} (zuò) — to do; to make}

\label{lesson:ZH-C89-zuo}



\begin{quote}
Before the new one: read the last lesson's character, then a word you already write.

\begin{itemize}
\raggedright
  \item \emph{Read it:} \zh{给你}, aloud — \emph{gěi nǐ}, here you are
  \item \emph{Recall:} say \zh{做饭} — \emph{zuòfàn}, to cook
  \item \emph{Recall:} say \emph{wǒmen zuò chē qù}, we're going by car
\end{itemize}
\end{quote}

\subsection*{You'll want to know: \zh{做}}
\textbf{\zh{做}} — \emph{zuò} — \textquotedblleft{}to do; to make\textquotedblright{}.

\begin{quote}
\textbf{\zh{你做什么}} \emph{nǐ zuò shénme} — What are you doing?
\end{quote}

\begin{quote}
\textbf{\zh{妈妈做饭}} \emph{māma zuò fàn} — Mum cooks.
\end{quote}

\begin{cousinweb}
\textbf{\zh{做}} (\emph{zuò}) has no reading in the Guangyun: no dictionary we use records an older form, so its history is not settled.
\end{cousinweb}

\begin{grammarlens}[title={What you've built}]
You have written \zh{做} since \zh{做饭}. \emph{zuò} 'to sit', as in \emph{zuò chē}, sounds the same and is a different word with a different character.
\end{grammarlens}

\subsection*{Your turn}
\begin{itemize}
\raggedright
  \item \emph{Say it:} \emph{zuò}
  \item \emph{Say it:} \emph{zuò}, once more
  \item \emph{Say it:} \zh{你做什么} — \emph{nǐ zuò shénme} — What are you doing?
  \item \emph{Say it:} \zh{妈妈做饭} — \emph{māma zuò fàn} — Mum cooks.
  \item \emph{Recall:} say \emph{zuò chē}, then \zh{做饭}, and hear the same sound twice
\end{itemize}

\begin{tcolorbox}[breakable,colback=teal!4,colframe=teal!35!black,title={Before you move on}]
What does \zh{做} mean? (\textquotedblleft{}To do; to make\textquotedblright{}.) Is it the \emph{zuò} of \emph{zuò chē}? (\textbf{No: same sound, another word}.)
\end{tcolorbox}
\section[fàng]{fàng — to put; to place}

\label{lesson:ZH-C89-fang}



\begin{quote}
Before the new one: last lesson's word, then a bag and a table.

\begin{itemize}
\raggedright
  \item \emph{Recall:} say \zh{你做什么} — \emph{nǐ zuò shénme}
  \item \emph{Recall:} say \emph{bāo}, a bag — \textbf{R3}, durable
  \item \emph{Recall:} say \emph{zhuōzi shàngbian}, on the table — \textbf{R2}
\end{itemize}
\end{quote}

\subsection*{You'll want to know: fàng}
\emph{fàng} — \textquotedblleft{}to put; to place\textquotedblright{}.

\begin{quote}
\emph{fàng zài zhèr} — Put it here.
\end{quote}

\begin{quote}
\emph{fàng zài zhuōzi shàngbian} — Put it on the table.
\end{quote}

Its character comes later; for now it is a word for the ear.

\begin{cousinweb}
\emph{fàng} is inherited from Middle Chinese \emph{pjangH} (Guangyun, in Baxter's transcription), glossed 'to drive out; to send away'. 'To put down' is a later sense the Guangyun does not give.
\end{cousinweb}

\begin{grammarlens}[title={What you've built}]
\emph{fàng}, then \zh{在}, then the place: the place words of the last two chapters all fit after it.
\end{grammarlens}

\subsection*{Your turn}
\begin{itemize}
\raggedright
  \item \emph{Say it:} \emph{fàng}
  \item \emph{Say it:} \emph{fàng}, once more
  \item \emph{Say it:} \emph{fàng zài zhèr} — Put it here.
  \item \emph{Say it:} \emph{fàng zài bāo lǐbian} — Put it in the bag.
  \item \emph{Recall:} say \emph{zuò}, then \emph{fàng} again
\end{itemize}

\begin{tcolorbox}[breakable,colback=teal!4,colframe=teal!35!black,title={Before you move on}]
What does \emph{fàng} mean? (\textquotedblleft{}To put; to place\textquotedblright{}.) Put it here. (\emph{fàng zài zhèr}.)
\end{tcolorbox}
\section[dài]{dài — to bring; to take along}

\label{lesson:ZH-C89-dai}



\begin{quote}
Before the new one: last lesson's word, then two things you take out with you.

\begin{itemize}
\raggedright
  \item \emph{Recall:} say \emph{fàng zài zhèr}, put it here — \textbf{R1}, one lesson back
  \item \emph{Recall:} say \zh{手机} — \emph{shǒujī}, a phone
  \item \emph{Recall:} say \emph{yàoshi}, a key
\end{itemize}
\end{quote}

\subsection*{You'll want to know: dài}
\emph{dài} — \textquotedblleft{}to bring; to take along\textquotedblright{}.

\begin{quote}
\emph{nǐ dài shǒujī le ma} — Have you brought your phone?
\end{quote}

Its character comes later; for now it is a word for the ear.

\begin{cousinweb}
\emph{dài} is inherited from Middle Chinese \emph{tajH} (Guangyun, in Baxter's transcription), glossed 'a belt; a sash'. How a belt became 'to carry along' is not given there, so the step is not settled.
\end{cousinweb}

\begin{grammarlens}[title={What you've built}]
\emph{nǐ dài shǒujī le ma}: \zh{了} after the verb says it is done, and \zh{吗} makes the question. The answer is \emph{dài le}, or \emph{méi dài}.
\end{grammarlens}

\subsection*{Your turn}
\begin{itemize}
\raggedright
  \item \emph{Say it:} \emph{dài}
  \item \emph{Say it:} \emph{dài}, once more
  \item \emph{Say it:} \emph{nǐ dài shǒujī le ma} — Have you brought your phone?
  \item \emph{Say it:} \emph{méi dài} — No, I haven't.
  \item \emph{Recall:} say \emph{yàoshi}, then \emph{dài} again
\end{itemize}

\begin{tcolorbox}[breakable,colback=teal!4,colframe=teal!35!black,title={Before you move on}]
What does \emph{dài} mean? (\textquotedblleft{}To bring; to take along\textquotedblright{}.) Say you didn't bring it. (\emph{méi dài}.)
\end{tcolorbox}
\section[zhàn]{zhàn — to stand}

\label{lesson:ZH-C89-zhan}



\begin{quote}
Before the new one: last lesson's word, then two places.

\begin{itemize}
\raggedright
  \item \emph{Recall:} say \emph{dài}, to bring — \textbf{R1}, one lesson back
  \item \emph{Recall:} say \emph{huǒchēzhàn}, a railway station — \textbf{R3}, durable
  \item \emph{Recall:} say \zh{门口} — \emph{ménkǒu}, the doorway
\end{itemize}
\end{quote}

\subsection*{You'll want to know: zhàn}
\emph{zhàn} — \textquotedblleft{}to stand\textquotedblright{}.

\begin{quote}
\emph{qǐng zhàn zài zhèr} — Please stand here.
\end{quote}

\begin{quote}
\emph{tā zhàn zài ménkǒu} — He's standing in the doorway.
\end{quote}

Its character comes later; for now it is a word for the ear.

\begin{cousinweb}
\emph{zhàn} is inherited from Middle Chinese \emph{treamH} (Guangyun, in Baxter's transcription), glossed 'in common speech, to stand alone'. The same word and character end \emph{huǒchēzhàn}, a railway station.
\end{cousinweb}

\begin{grammarlens}[title={What you've built}]
\emph{zhàn} takes \zh{在} and a place, as \emph{fàng} does: \emph{zhàn zài ménkǒu}, standing in the doorway.
\end{grammarlens}

\subsection*{Your turn}
\begin{itemize}
\raggedright
  \item \emph{Say it:} \emph{zhàn}
  \item \emph{Say it:} \emph{zhàn}, once more
  \item \emph{Say it:} \emph{qǐng zhàn zài zhèr} — Please stand here.
  \item \emph{Say it:} \emph{tā zhàn zài ménkǒu} — He's standing in the doorway.
  \item \emph{Recall:} say \emph{huǒchēzhàn}, then \emph{zhàn} again
\end{itemize}

\begin{tcolorbox}[breakable,colback=teal!4,colframe=teal!35!black,title={Before you move on}]
What does \emph{zhàn} mean? (\textquotedblleft{}To stand\textquotedblright{}.) Where else have you heard it? (\textbf{At the end of \emph{huǒchēzhàn}}.)
\end{tcolorbox}
\section[Practice]{\zh{给我}}

\label{lesson:ZH-C89-practice}



\begin{quote}
The chapter closes here, with somebody getting out of the door.

\begin{itemize}
\raggedright
  \item \emph{Recall:} say \emph{zhàn}, to stand — \textbf{R1}, one lesson back
  \item \emph{Recall:} say \emph{zhǔnbèi}, to get ready — \textbf{R3}, durable
  \item \emph{Recall:} say \emph{chūmén}, to go out — \textbf{R4}, recognition at distance
\end{itemize}
\end{quote}

\subsection*{The exchange}
\begin{quote}
\textbf{A:} \zh{你做什么}
\end{quote}

\begin{quote}
\emph{nǐ zuò shénme} — What are you doing?
\end{quote}

\begin{quote}
\textbf{B:} \emph{wǒ zhǔnbèi chūmén} — I'm getting ready to go out.
\end{quote}

\begin{quote}
\textbf{A:} \emph{nǐ dài yàoshi le ma} — Have you got your key?
\end{quote}

\begin{quote}
\textbf{B:} \emph{dài le. shǒujī ne} — Yes. Where's my phone?
\end{quote}

\begin{quote}
\textbf{A:} \emph{zài zhuōzi shàngbian. gěi nǐ} — On the table. Here you are.
\end{quote}

\begin{quote}
\textbf{B:} \emph{xièxie. fàng zài bāo lǐbian} — Thanks. In the bag it goes.
\end{quote}

\begin{quote}
\textbf{A:} \emph{qǐng zhàn zài ménkǒu, wǒ kànkan} — Stand in the doorway, let me look at you.
\end{quote}

Five actions in one doorway: do, bring, give, put and stand. \zh{给你}, here you are, is \zh{给} handing something over as it is said.

\begin{tcolorbox}[breakable,skin=enhanced,colback=orange!6,colframe=orange!45!black,boxrule=0.5pt,arc=1mm,left=6pt,right=6pt,top=4pt,bottom=4pt,fonttitle=\bfseries,title={Read this}]
\emph{Read it:} the exchange through once without stopping, then answer the question below
\end{tcolorbox}

\begin{grammarlens}[title={What you've built — five actions}]
Five verbs, and the three that take a place take it with \zh{在}:

\begin{itemize}
\raggedright
  \item \textbf{\zh{给我}} \emph{gěi wǒ} — give it to me
  \item \textbf{\zh{你做什么}} \emph{nǐ zuò shénme} — what are you doing?
  \item \emph{fàng zài zhèr} — put it here
  \item \emph{nǐ dài shǒujī le ma} — have you brought your phone?
  \item \emph{qǐng zhàn zài zhèr} — please stand here
\end{itemize}

\textbf{Listen.} \emph{Gěi wǒ nǐ de shǒujī.}

\textbf{Speak.} Hand it over with \zh{给你}, then say where you put it.

\emph{Write it:} \zh{给} — the silk piece first, then the roof, the bar and \zh{口}
\end{grammarlens}

\subsection*{Your turn}
\begin{itemize}
\raggedright
  \item \emph{Say it:} \zh{给我} — \emph{gěi wǒ} — Give it to me.
  \item \emph{Say it:} \zh{给你} — \emph{gěi nǐ} — Here you are.
  \item \emph{Say it:} \emph{nǐ dài yàoshi le ma} — Have you got your key?
  \item \emph{Say it:} the whole exchange, both parts
  \item \emph{Write it:} \zh{给你}
\end{itemize}

\begin{tcolorbox}[breakable,colback=teal!4,colframe=teal!35!black,title={Before you move on}]
Hand something over. (\textbf{\zh{给你}}.) Ask what somebody is doing. (\textbf{\zh{你做什么}}.) Tell somebody to put it here. (\emph{fàng zài zhèr}.) Ask whether they brought their key. (\emph{nǐ dài yàoshi le ma}.)
\end{tcolorbox}
